% ============================================================================
\chapter{网卡驱动}
\label{ch:nic-driver}
% Stage H-1 ~ H-3
% ============================================================================

\section{本章目标}

\begin{itemize}
  \item 理解网络分层模型（OSI / TCP-IP）与以太网帧结构
  \item 设计可插拔 \texttt{netdev\_ops\_t} 接口，解耦网卡驱动
  \item 实现 PCI 设备枚举与 MMIO 映射
  \item 实现 Intel E1000 驱动后端（QEMU 默认网卡）
  \item 了解 virtio-net 备选方案
  \item 通过 \texttt{KCONFIG\_NETDEV\_BACKEND} 一键切换
\end{itemize}

% ============================================================================
\section{概念：网络分层}
\label{sec:nic-concept}
% ============================================================================

% TODO:
% - TCP/IP 四层模型（从下到上）：
%     1. 链路层（Link）：以太网帧收发、MAC 地址
%     2. 网络层（Network）：IP 路由、ICMP
%     3. 传输层（Transport）：TCP/UDP
%     4. 应用层（Application）：HTTP、DNS
% - 以太网帧结构：
%     [6B dst MAC][6B src MAC][2B EtherType][46~1500B payload][4B FCS]
%     EtherType: 0x0800=IPv4, 0x0806=ARP
% - MTU = 1500 bytes（不含帧头帧尾）
% - 图：数据封装过程（HTTP → TCP 段 → IP 包 → 以太网帧）

% ============================================================================
\section{可插拔接口：\texttt{netdev\_ops\_t}}
\label{sec:nic-interface}
% ============================================================================

% TODO:
% - 接口设计：
%     typedef struct netdev_ops {
%         const char* name;                          // "e1000", "virtio_net"
%         int (*init)(void);                          // 初始化网卡硬件
%         int (*send)(const void* frame, size_t len); // 发送以太网帧
%         int (*recv)(void* buf, size_t max_len);     // 接收以太网帧
%         void (*get_mac)(uint8_t mac[6]);            // 获取 MAC 地址
%         int (*link_status)(void);                   // 是否连接（1/0）
%     } netdev_ops_t;
% - dispatch 层 + kconfig.h: KCONFIG_NETDEV_BACKEND 0=e1000, 1=virtio_net

% ============================================================================
\section{PCI 设备枚举}
\label{sec:nic-pci}
% ============================================================================

% TODO:
% - PCI 配置空间读取（I/O ports 0xCF8/0xCFC）
% - 遍历 bus/device/function → 读 vendor_id/device_id
% - E1000 标识：vendor=0x8086, device=0x100E (82540EM)
% - BAR0 → MMIO 基地址 → vmm_map 映射到内核虚拟地址

% ============================================================================
\section{后端一：E1000 驱动}
\label{sec:nic-e1000}
% ============================================================================

% TODO:
% - E1000 初始化序列：
%     1. 读取 MMIO BAR0 基地址 → 映射到虚拟地址
%     2. 软复位（CTRL.RST）
%     3. 读取 MAC 地址（RAL/RAH 寄存器或 EEPROM）
%     4. 配置 RX：分配 RX 描述符环 + 缓冲区（DMA 物理地址）→ 写 RDBAL/RDBAH/RDLEN/RDH/RDT
%     5. 配置 TX：分配 TX 描述符环 → 写 TDBAL/TDBAH/TDLEN/TDH/TDT
%     6. 启用 RX/TX（RCTL.EN + TCTL.EN）
%     7. 配置中断（IMS 寄存器）+ 注册 IRQ handler
% - 发送：构造以太网帧 → 填 TX 描述符 → 更新 TDT → 等待完成
% - 接收：IRQ 通知有新帧 → 从 RX 描述符读取 → 传递给协议栈 → 更新 RDT
% - [WHY] TX/RX 使用描述符环：DMA 让网卡直接读写物理内存，CPU 不需逐字节搬运

% ============================================================================
\section{后端二：virtio-net（备选）}
\label{sec:nic-virtio}
% ============================================================================

% TODO:（备选，可跳过）
% - virtio 模型：Guest 和 Host 通过共享内存 virtqueue 通信
% - 比 E1000 仿真更高效（少很多寄存器操作）
% - PCI device: vendor=0x1AF4, device=0x1000
% - KCONFIG_NETDEV_BACKEND=1 切换

% ============================================================================
\section{验证}
\label{sec:nic-verify}
% ============================================================================

% TODO:
% - 测试 1：PCI 枚举 → 找到 E1000 → 输出 vendor/device/BAR
% - 测试 2：读取并打印 MAC 地址
% - 测试 3：发送一个 ARP 请求帧 → Wireshark/tcpdump 在 host 抓到
% - 测试 4：接收来自 host 的帧 → 解析并打印
% - shell 命令：net info / net send / net recv

% ============================================================================
\section{本章小结}
% ============================================================================

网卡驱动通过可插拔 \texttt{netdev\_ops\_t} 接口，让 E1000 和 virtio-net 可以一键切换。
E1000 的 TX/RX 描述符环 + DMA 是理解现代网卡驱动的基础。
下一章在驱动之上构建 TCP/IP 协议栈，让 OS 能进行真正的网络通信。
